\documentclass[3p,sort&compress,11pt,fleqn,review]{elsarticle}
\bibliographystyle{elsarticle-num-names}
\usepackage{amsmath,amsfonts,amssymb,siunitx,float,subcaption,setspace,booktabs,bm,mathtools,url,algorithm,algpseudocode,lineno,lmodern}
\usepackage[fixed]{fontawesome5}
%\usepackage{easyReview}
\journal{}
\newcommand{\balpha}{\mb{\alpha}}
\newcommand{\bbeta}{\mb{\beta}}
\newcommand{\bc}{\mb{c}}
\newcommand{\bd}{\mb{d}}
\newcommand{\be}{\mb{e}}
\newcommand{\beeta}{\mb{\eta}}
\newcommand{\bgamma}{\mb{\gamma}}
\newcommand{\bn}{\mb{n}}
\newcommand{\bq}{\mb{q}}
\newcommand{\bs}{\mb{s}}
\newcommand{\bsigma}{\mb{\sigma}}
\newcommand{\bv}{\mb{v}}
\newcommand{\bvarepsilon}{\mb{\varepsilon}}
\newcommand{\bzeta}{\mb{\zeta}}
\newcommand*{\algoref}[1]{Algorithm~\ref{#1}}
\newcommand*{\ddfrac}[2]{\dfrac{\md{#1}}{\md{#2}}}
\newcommand*{\dev}[1]{\text{dev}\left(#1\right)}
\newcommand*{\diag}[1]{\text{diag}\left(#1\right)}
\newcommand*{\eqsref}[1]{Eq.~(\ref{#1})}
\newcommand*{\figref}[1]{Fig.~\ref{#1}}
\newcommand*{\mb}[1]{\bm{#1}}
\newcommand*{\md}[1]{\mathrm{d}#1}
\newcommand*{\mdd}[1]{\mathrm{d}^2#1}
\newcommand*{\mT}{\mathrm{T}}
\newcommand*{\pdfrac}[2]{\dfrac{\partial{#1}}{\partial{#2}}}
\newcommand*{\secref}[1]{\S~\ref{#1}}
\newcommand*{\sign}[1]{\text{sign}\left(#1\right)}
\newcommand*{\tabref}[1]{Table~\ref{#1}}
\DeclarePairedDelimiter\abs{\lvert}{\rvert}
\DeclarePairedDelimiter\norm{\lVert}{\rVert}
\makeatletter
\let\oldabs\abs\def\abs{\@ifstar{\oldabs}{\oldabs*}}
\let\oldnorm\norm\def\norm{\@ifstar{\oldnorm}{\oldnorm*}}
\newenvironment{breakablealgorithm}{\begin{center}\refstepcounter{algorithm}\hrule height.8pt depth0pt \kern2pt
\renewcommand{\caption}[2][\relax]{
{\raggedright\textbf{\ALG@name~\thealgorithm} ##2\par}
\ifx\relax##1\relax
\addcontentsline{loa}{algorithm}{\protect\numberline{\thealgorithm}##2}
\else
\addcontentsline{loa}{algorithm}{\protect\numberline{\thealgorithm}##1}
\fi
\kern2pt\hrule\kern2pt}}{\kern2pt\hrule\relax\end{center}}
\makeatother
\begin{document}
\linenumbers
\begin{abstract}
    \begin{linenumbers}
        A versatile constitutive model is proposed.
    \end{linenumbers}
\end{abstract}
\begin{keyword}
    constitutive modelling\sep
    plasticity\sep
    metal\sep
    ratcheting\sep
    cyclic response
\end{keyword}
\begin{frontmatter}
    \title{The \texttt{Balloon}-v1 Model}
    \author[add1]{Theodore~L.~Chang\corref{tlc}}\ead{tlcfem@gmail.com}
    \cortext[tlc]{corresponding author}
    \address[add1]{IRIS Adlershof, Humboldt-Universität zu Berlin, Berlin, Germany, 12489.}
\end{frontmatter}
\section{Background}
\subsection{The Basic Framework}
\subsubsection{Constitutive Law}
Based on the additive decomposition of deformation, the conventional elasticity applies.
That is, the stress tensor $\bsigma$ can be expressed as the product of the elasticity tensor $\mb{D}_e$ and the elastic strain tensor $\bvarepsilon^e$,
\begin{gather}
    \bsigma=\mb{D}_e:\bvarepsilon^e=\mb{D}_e:\left(\bvarepsilon-\bvarepsilon^p\right),
\end{gather}
where $\bvarepsilon$ is the total strain tensor and $\bvarepsilon^p$ is the plastic strain tensor.
\subsubsection{Subloading Surface}
The subloading surface differs from the conventional yield surface by an additional scalar field $z$, which is named as the normal yield ratio in the existing literature.
Adopting the von Mises yield criterion, the equivalent subloading surface \citep{Chang2025} can be expressed as
\begin{gather}
    f_s=\norm{\beeta}-z\sqrt{\dfrac{2}{3}}\sigma^y,
\end{gather}
where $\beeta=\dev{\bsigma}-a^y\balpha+\left(z-1\right)\sigma^y\bd$ is the shifted stress tensor,
$\balpha$ is the conventional back stress tensor, $\bd$ is the similarity tensor, both are internal history fields that would evolve in the stress space. The scalars $a^y$ and $\sigma^y$ are the magnitudes that may also evolve with the development of plasticity.
\subsubsection{Flow Rule}
Assuming associated plasticity, the flow rule possesses the widely known form,
\begin{gather}
    \dot{\bvarepsilon^p}=\gamma{}\bn,
\end{gather}
where $\gamma$ is the plastic multiplier, and $\bn$ is the unit gradient of $f_s$, that is,
\begin{gather}
    \bn=\pdfrac{f_s}{\bsigma}/\norm{\pdfrac{f_s}{\bsigma}}.
\end{gather}
\subsubsection{Evolution Rules}
The scalar fields are driven by the accumulated plastic strain $q$, its rate form is defined as
\begin{gather}
    \dot{q}=\sqrt{\dfrac{2}{3}}\gamma.
\end{gather}
With the above, the following evolution rule is introduced for $z$,
\begin{gather}
    \dot{z}=U\dot{q},
\end{gather}
where $U=U\left(z\right)$ is a function of $z$.
Furthermore,
\begin{gather}
    \sigma^y=\sigma^i+k_\text{iso}q+\sigma^s_\text{iso}\left(1-e^{-m^s_\text{iso}q}\right),\\
    a^y=a^i+k_\text{kin}q+a^s_\text{kin}\left(1-e^{-m^s_\text{kin}q}\right).
\end{gather}
The internal history fields (tensors) adopt the following evolution rules,
\begin{gather}
    \dot{\balpha}=b\left(\sqrt{\dfrac{2}{3}}\bn-\balpha\right)\gamma,\qquad
    \dot{\bd}=c_e\left(\sqrt{\dfrac{2}{3}}z_e\bn-\bd\right)\gamma.
\end{gather}
All unmentioned symbols/parameters are material constants.
\section{Limitations}
The most distinctive feature of the subloading surface model is that the normal yield ratio $z$ gradually approaches unity from zero during the loading process.
Thus, there is no elastic region during this process, as $z$ is driven by $q$, which accumulates with the plastic deformation.
\section{Towards A Better One}
\subsection{History of Load Reversals}
In conventional plasticity models, local reversals cannot be easily detected.
A common criterion for detecting reversal is the double contraction $\beeta:\dot{\bs}$.
For example, in the multi-surface theory \citep{Semenov2022}, the exact criterion is used to identify the initiation of a load reversal, which further controls the activation of a new surface.

The normal yield ratio $z$ effectively represents the current `loading level' at any given moment.
When it increases, the subloading surface expands, indicating a loading process.
When it decreases, the subloading surface contracts, indicating an unloading process.
Thus, with the introduction of $z$, local reversals can be detected with ease by checking whether $z$ ever decreases during the given (sub)step.

Let the normal yield ratio at each load reversal be denoted as $z_r^j$ with $j=1,2,3,\cdots$ standing for the $j$-th reversal.
The initial value is $z_r^0=0$.
We use $z_r$ to denote the collection of all $z_r^j$, thus,
\begin{gather}
    z_r=\{z_r^j\}.
\end{gather}
\section{Numerical Demonstration}
\subsection{Isotropic Hardening Stagnation}
\subsection{Cyclic Hardening}
\subsection{Cyclic Softening}
\subsection{Improved Ratcheting}
\bibliography{BIB}
\end{document}